\chapter{Resultados}

En este apartado se presentan los resultados derivados del proceso de implementación del videojuego Ascension, describiendo el estado final del sistema, las funcionalidades alcanzadas y las limitaciones técnicas identificadas tras la ejecución de pruebas funcionales.

\section{Sistema final implementado}

Se ha implementado un videojuego completo del subgénero \textit{rogue-like} bidimensional con vista cenital, que abarca un ciclo jugable completo desde el menú principal hasta la condición de victoria o derrota. El sistema integra un conjunto cohesionado de subsistemas que colaboran para generar una experiencia de juego funcional y extensible.

La arquitectura del sistema se ha construido mediante una separación clara de responsabilidades entre componentes. El gestor global de estados (\textit{GameManager}) coordina la transición entre escenas y mantiene la coherencia del flujo de juego. La gestión de salas (\textit{RoomFlowController}) controla el avance entre niveles y la aparición de enemigos según criterios predefinidos. El generador de salas (\textit{RoomGenerator}) produce de forma procedural la estructura espacial de cada nivel mediante algoritmos basados en \textit{tilemaps}.

El proyecto se estructura en tres escenas principales: \textit{MenuPrincipal}, \textit{SeleccionDeClase} y \textit{GameScene}. La escena \textit{GameScene} constituye el núcleo del juego, donde transcurre la totalidad del flujo mediante la regeneración dinámica de salas dentro de la misma escena.

\section{Generación procedural de salas y flujo de juego}

El flujo de juego avanza mediante un sistema de salas lineales controlado por un contador interno. El jugador accede a la siguiente sala mediante la interacción con escaleras de transición, que se activan una vez que la sala actual ha sido limpiada de enemigos. Cada nueva sala incrementa dicho contador, el cual determina la aparición periódica de encuentros especiales cada cinco salas.

La generación procedural se ha validado verificando que las salas generadas presentan muros de contención íntegros, un área jugable accesible sin obstáculos, puntos de transición correctamente posicionados y que la regeneración del suelo no interfiere con las mecánicas de colisión de enemigos y jugador.

\section{Sistema de combate, armas y enemigos}

Se ha implementado un sistema de combate que distingue entre ataques cuerpo a cuerpo y ataques a distancia. Las armas cuerpo a cuerpo realizan un escaneo de área alrededor del jugador en la dirección de orientación, aplicando daño a todos los enemigos presentes dentro del rango de ataque. Las armas a distancia generan proyectiles que se desplazan en línea recta y aplican daño al entrar en colisión con enemigos.

El juego incluye tres tipos de enemigos base con comportamientos diferenciados: enemigos perseguidores que se desplazan directamente hacia el jugador, enemigos agresivos que atacan cuando se encuentran a una distancia reducida, y enemigos de rango que mantienen la distancia y realizan ataques a distancia. Cada tipo de enemigo se configura mediante datos almacenados en \textit{ScriptableObject}, lo que permite parametrizar atributos como velocidad, daño, vida y radio de visión sin modificar la lógica de ejecución.

La aparición de enemigos se gestiona mediante un sistema de generación ponderada por coste (\textit{EnemyManager}), que instancia enemigos de forma aleatoria dentro de un radio máximo de seguridad respecto al jugador. Este mecanismo evita la aparición de enemigos en la proximidad inmediata del jugador, mejorando la experiencia de juego.

El sistema de combate se ha validado verificando que los ataques cuerpo a cuerpo aplican daño correctamente a los enemigos dentro del rango de escaneo, que los proyectiles siguen la trayectoria calculada y aplican daño al impactar, que los enemigos se desplazan conforme a su comportamiento definido, y que la eliminación de enemigos se refleja correctamente en el control de salas limpias.

\section{Sistema de jefes y condición de victoria}

Se ha implementado un subsistema de enemigos especiales que aparecen de forma periódica cada cinco salas completadas. Estos enemigos no introducen nuevas mecánicas de inteligencia artificial respecto a los enemigos base, sino que consisten en variantes escaladas de los enemigos normales, con incrementos en atributos como vida total y daño infligido.

La aparición de estos enemigos se produce en posiciones alejadas del jugador, calculadas mediante puntos predefinidos dentro de la sala, con el objetivo de garantizar una entrada progresiva al combate y evitar situaciones de daño inmediato injustificado.

La condición de victoria se define como la superación de tres encuentros de este tipo. Cada enemigo especial derrotado incrementa un contador interno que, al alcanzar dicho valor, desencadena el evento de victoria del juego. Al cumplirse esta condición se muestra la pantalla de victoria, desde la cual es posible retornar al menú principal o iniciar una nueva partida.

La validación de este subsistema ha confirmado que el contador de salas controla correctamente la aparición de estos encuentros, que las estadísticas de los enemigos especiales se escalan adecuadamente respecto a los enemigos base y que el evento de victoria se activa conforme a lo previsto.

\section{Estabilidad y coherencia del sistema}

Tras la ejecución de múltiples sesiones de pruebas funcionales se ha validado la estabilidad del flujo completo de juego. Se ha verificado que la transición entre escenas no genera errores de referencia nula, que el \textit{GameManager} mantiene la coherencia del estado global durante toda la ejecución, que la generación procedural de salas no produce inconsistencias geométricas, que los enemigos se destruyen correctamente al cambiar de sala y que la persistencia de puntuación mediante \textit{PlayerPrefs} funciona sin corrupción de datos.

El sistema es capaz de completar múltiples ciclos completos de juego, desde el menú principal hasta la condición de victoria, sin presentar fallos críticos ni pérdida de información. La interfaz de usuario responde adecuadamente a los cambios de estado y muestra información actualizada del jugador durante todo el flujo de juego.

No se han implementado pruebas automatizadas con cobertura estadística, por lo que la validación se ha centrado en la detección de errores durante ejecuciones manuales en el editor de Unity.

\section{Herramientas de automatización desarrolladas}

De forma paralela al desarrollo del videojuego, se han implementado herramientas específicas dentro del editor de Unity para automatizar tareas de configuración y validación. Estas herramientas incluyen un generador automático de estructuras de escena con posicionamiento de elementos de interfaz de usuario, un editor visual de parámetros de enemigos con validaciones en tiempo de edición, un sistema de aplicación masiva de efectos de baldosas mediante operaciones por lotes, y utilidades de depuración que visualizan información de estado durante la ejecución.

El uso de estas herramientas ha reducido notablemente la incidencia de errores de configuración y ha permitido iterar de forma más ágil sobre parámetros de balance sin necesidad de recompilar código.

\section{Limitaciones técnicas y aspectos no implementados}

Se han identificado las siguientes limitaciones técnicas y aspectos no implementados por restricciones de tiempo y alcance:

\begin{itemize}
    \item El sistema de recogida de armas no integra completamente la lógica de equipamiento automático en el controlador del jugador (\textit{PlayerController}), requiriendo integración adicional para su funcionamiento completo
    \item No se ha implementado un sistema integral de audio, incluyendo gestión de música de fondo, efectos de sonido y transiciones según el estado del juego
    \item El sistema de caídas de objetos se limita a armas y efectos de baldosas, sin incluir un sistema genérico de probabilidades para la generación de consumibles u otros tipos de ítems
    \item La inteligencia artificial de enemigos implementa comportamientos básicos de persecución y ataque, sin incorporar tácticas avanzadas como coordinación entre enemigos o patrones complejos de ataque
    \item No se ha realizado una fase exhaustiva de balance de progresión de armas, pesos de aparición de enemigos y escalado de dificultad. Los parámetros actuales permiten completar el juego, pero un proceso de balance con pruebas de usuario habría mejorado la experiencia final
    \item No se ha implementado un sistema de persistencia de progreso entre partidas, salvo la puntuación, por lo que cada ejecución comienza desde un estado inicial sin desbloqueables ni mejoras permanentes
    \item Finalmente, la documentación interna del código se ha realizado de forma parcial mediante comentarios en métodos críticos, sin alcanzar una cobertura exhaustiva de todas las clases y métodos del sistema.
\end{itemize}

En conjunto, los resultados obtenidos confirman que el sistema desarrollado cumple los objetivos técnicos establecidos para el proyecto, proporcionando una base funcional, coherente y extensible para futuras mejoras.

Los resultados obtenidos han permitido validar el correcto funcionamiento del sistema en condiciones normales de ejecución. No obstante, algunas funcionalidades previstas, como la incorporación de música, un sistema avanzado de generación de armas raras y un proceso exhaustivo de balanceo, no han sido implementadas por limitaciones temporales. Estas carencias no comprometen el núcleo funcional del sistema ni la consecución de los objetivos técnicos planteados.
